\documentclass[11pt,letterpaper]{article}
\usepackage[top=0.6in, bottom=0.45in, left=0.65in, right=0.65in]{geometry}
\usepackage{lmodern}
\usepackage{hyperref}
\hypersetup{colorlinks=false, urlcolor=black, pdfborder=0 0 0}
\setlength{\parindent}{0pt}
\setlength{\parskip}{0.3em}
\raggedright
\pagestyle{empty}

\begin{document}

\noindent
{\bfseries Ke-Cheng Hsu (Clio)}\\
Ph.D.\ Candidate, Department of Economics, HKUST\\
(852) 6840\,9753 \quad $|$ \quad \href{mailto:khsu@connect.ust.hk}{khsu@connect.ust.hk}\\
\href{https://cliohsu1997.github.io/}{cliohsu1997.github.io} \quad $|$ \quad \href{https://github.com/cliohsu1997}{github.com/cliohsu1997}

\vspace{0.4em}
\noindent
November 12, 2026

\vspace{0.35em}
\noindent
Hiring Team\\
UBS --- Research, European Equity and Derivatives Strategy, Hyderabad\\
Re: Research -- Economics \& Strategy, Equity \& Derivative Associate (345504BR)

\vspace{0.2em}
Dear Hiring Team,

I am writing to apply for the Research -- Economics \& Strategy, Equity \& Derivative Associate position in Hyderabad (345504BR). I am a Ph.D.\ candidate in Economics at HKUST, working on empirical industrial organization and digital economics, and I expect to graduate in June 2027 (I can finish earlier if that helps). Your team maintains the analytical frameworks behind regional, country, sector, and style views and researches the macro developments, market dynamics, and investor behavior that drive equity, volatility, and derivative markets. That is the kind of work I am trained for: build a framework that ties evidence to theory, estimate it carefully on large data, and write up the result so a non-specialist can follow it.

The day-to-day described in the posting is close to what I already do in research. In my job market paper on Taobao apparel return policy (2017--2023; \(\sim\)62M RMB revenue; six categories; \(\sim\)20--30\% return rates), I clean large transaction data, estimate a structural model of purchase, return-shipping-insurance take-up, and return decisions, and rank policies by seller profit, consumer surplus, and social welfare under what-if scenarios: selective adoption outperforms automatic provision by \(\sim\)15--30 pp of first-best welfare in thinner-margin categories, and a designed partial return-penalty discount attains \(\sim\)75--78\% of first-best welfare. Maintaining a country, sector, or style view is the same exercise: an explicit framework, estimated on data, that says what changes when conditions change.

With Professors Kohei Kawaguchi and Kei Kawai, I build title-by-market-by-month panels from monthly bookstore point-of-sale records (tens of billions of records across 1{,}000+ stores) and estimate pricing and inventory under demand uncertainty. We show that the classical static model overstates marginal cost by \(\sim\)20\% --- a correction that changes a company's assessed value. Findings go out as figures, tables, and a weekly report on a fixed deadline, and I have presented this work at the Asia-Pacific Industrial Organization Conference, the Econometric Society Australasia Meeting, and the China Economics Annual Conference in 2026 --- the same cadence as the team's investment themes, publications, and thought leadership.

On the tools side, I run tested, reproducible pipelines in Python, R, and C++ on AWS and HPC/SLURM, and I use Microsoft Excel for figures and summary tables. I use AI tools every day for research, coding, and OCR data extraction, with outputs verified by tests, so the AI-assisted research and workflow automation the posting asks for is already my standard workflow. My earlier research modeled a bank crisis on an interbank lending network in MATLAB, so financial-market dynamics and systemic-risk questions are familiar territory.

Explaining the numbers and working with people is the other half of this job. As Teaching Assistant Coordinator for the Business School at HKUST (2025--2026), I coached and evaluated a team of teaching assistants and held their teaching materials and mock lectures to one standard --- the quality control a research team applies before a note goes out. As Hall Tutor at UG Hall I (2024--present), I mediate resident conflicts and lead orientation, workshops, and career talks. My teaching is rated 4.53/5 as a lecturer and 4.72/5 on average as a teaching assistant --- years of practice making technical material clear to non-specialists. I am happy to be based in Hyderabad, working with strategists in London and the US across time zones.

My resume is attached, and I can send my job market paper on request. Thank you for your consideration; I hope we get the chance to talk.

\vspace{0.4em}
Sincerely,\\[0.3em]
Ke-Cheng Hsu (Clio)

\end{document}
